% =====================================================================
%  LDA-T318 — Large Language Models and Generative AI for NLP
%  Shared Beamer preamble for all lecture decks.
%  Aarne Talman, University of Helsinki, 2026.
%
%  Engine:  XeLaTeX or LuaLaTeX preferred (gives the Fira font set).
%           pdfLaTeX also works; Metropolis falls back to Latin Modern.
%  Build:   latexmk -xelatex main.tex      (from the deck's own directory)
%
%  Every deck does:
%      \documentclass[aspectratio=169,10pt]{beamer}
%      \input{../common/preamble.tex}
% =====================================================================

% ---------------------------------------------------------------- theme
\usetheme{metropolis}
\metroset{
  titleformat=smallcaps,
  numbering=fraction,
  progressbar=frametitle,
  block=fill,
  sectionpage=progressbar,
  subsectionpage=none,
}

% ---------------------------------------------------------------- packages
% inputenc/fontenc are required under pdfLaTeX and actively harmful under
% XeLaTeX/LuaLaTeX, where Metropolis loads fontspec to get Fira. Guard them.
\usepackage{iftex}
\ifPDFTeX
  \usepackage[utf8]{inputenc}
  \usepackage[T1]{fontenc}
  \usepackage[sfdefault]{FiraSans}
  \usepackage{FiraMono}
\else
  % Metropolis only gets Fira if fontspec can actually find it. TeX Live ships
  % the OTFs but does NOT register them with fontconfig, so a plain
  % \setsansfont{Fira Sans} silently falls back to Latin Modern and the deck
  % stops looking like Metropolis. Load them by filename instead.
  \usepackage{fontspec}
  \setsansfont{FiraSans}[
    Extension      = .otf,
    UprightFont    = *-Book,
    ItalicFont     = *-BookItalic,
    BoldFont       = *-Medium,
    BoldItalicFont = *-MediumItalic]
  \setmonofont{FiraMono}[
    Extension   = .otf,
    UprightFont = *-Regular,
    BoldFont    = *-Medium,
    Scale       = 0.92]
  \renewcommand{\familydefault}{\sfdefault}
\fi
\usepackage[english]{babel}

\usepackage{amsmath,amssymb}
\usepackage{booktabs}
\usepackage{array}
\usepackage{multicol}
\usepackage{xcolor}
\usepackage{graphicx}
\usepackage{tikz}
\usepackage{pgfplots}
\usepackage{listings}
\usepackage{appendixnumberbeamer}

\pgfplotsset{compat=1.17}
\usetikzlibrary{positioning,arrows.meta,fit,backgrounds,calc,matrix,shapes.geometric,decorations.pathreplacing}

% =====================================================================
%  University of Helsinki palette
%  Primary brand colour is UH Hray/black with the "HY blue" accent.
%  Secondary accents chosen for projector contrast, not for print.
% =====================================================================
\definecolor{UHBlue}{HTML}{0C5A9E}   % primary accent
\definecolor{UHDark}{HTML}{15202B}   % body text / dark backgrounds
\definecolor{UHGrey}{HTML}{5A6672}   % secondary text
\definecolor{UHLight}{HTML}{EEF2F6}  % block backgrounds
\definecolor{UHOrange}{HTML}{D1601A} % alerts, gotchas
\definecolor{UHGreen}{HTML}{1F7A5C}  % examples, "this works"
\definecolor{UHRed}{HTML}{A62B2B}    % failure modes
\definecolor{UHSand}{HTML}{C8B88A}   % neutral third accent

\setbeamercolor{normal text}{fg=UHDark, bg=white}
\setbeamercolor{alerted text}{fg=UHOrange}
\setbeamercolor{example text}{fg=UHGreen}
\setbeamercolor{frametitle}{fg=white, bg=UHBlue}
\setbeamercolor{title separator}{fg=UHBlue}
\setbeamercolor{progress bar}{fg=UHBlue}
\setbeamercolor{progress bar in head/foot}{fg=UHBlue}
\setbeamercolor{progress bar in section page}{fg=UHBlue}
\setbeamercolor{structure}{fg=UHBlue}
\setbeamercolor{palette primary}{fg=white, bg=UHBlue}   % section pages, standout frames
\setbeamercolor{block title}{fg=white, bg=UHBlue}
\setbeamercolor{block body}{fg=UHDark, bg=UHLight}
\setbeamercolor{block title alerted}{fg=white, bg=UHOrange}
\setbeamercolor{block body alerted}{fg=UHDark, bg=UHLight}
\setbeamercolor{block title example}{fg=white, bg=UHGreen}
\setbeamercolor{block body example}{fg=UHDark, bg=UHLight}
\setbeamercolor{footnote}{fg=UHGrey}

\setbeamerfont{frametitle}{size=\large, series=\bfseries}
\setbeamerfont{footnote}{size=\tiny}
\setbeamerfont{caption}{size=\scriptsize}

% Frame title bar.
%
% Do NOT reimplement the frametitle template here. Metropolis's own template
% already spans \paperwidth (so `bg=UHBlue' above gives the full-width bar),
% already struts two-line titles without collision, and already draws the
% frametitle progress bar. Reimplementing it means touching internals whose
% names contain `@', which is NOT a letter in a .tex file -- `\@empty' parses
% as `\@' followed by the word "empty" and the build breaks.
%
% To change the bar's height, adjust Metropolis's own padding length instead,
% inside \makeatletter so that `@' is a letter:
%
%   \makeatletter
%   \setlength{\metropolis@frametitle@padding}{2.4ex}  % default 2.2ex
%   \makeatother
%
\makeatletter
\setlength{\metropolis@frametitle@padding}{1.2ex}
\makeatother

% No deck uses \framesubtitle; if one ever does, Metropolis ignores it and the
% subtitle should go in the frame body instead.

% =====================================================================
%  Course metadata — decks override \lecturename and \lecturenumber
% =====================================================================
\newcommand{\coursecode}{LDA-T318}
\newcommand{\coursename}{LLMs and Generative AI for NLP}
\newcommand{\courseteacher}{Aarne Talman}
\newcommand{\courseinstitution}{University of Helsinki}
\providecommand{\lecturename}{}
\providecommand{\lecturedate}{}

% Footline: course code on the left, lecture name centred, n/N on the right.
\setbeamertemplate{frame footer}{\coursecode\quad\textbar\quad\lecturename}

% =====================================================================
%  Semantic macros — use these instead of ad-hoc formatting, so that a
%  change of house style is a single edit here.
% =====================================================================

% --- Sourcing. Every quantitative claim on a slide should carry one. ---
\newcommand{\source}[1]{%
  \vfill{\scriptsize\color{UHGrey}\raggedright Source: #1\par}%
}
\newcommand{\sources}[1]{%
  \vfill{\scriptsize\color{UHGrey}\raggedright Sources: #1\par}%
}
% Inline citation chip, e.g. \cit{Raschka 2026a}
\newcommand{\cit}[1]{{\scriptsize\color{UHGrey}[#1]}}

% --- Counter-intuitive point. Three per lecture; these are exam material. ---
\newenvironment{gotcha}[1][Gotcha]{%
  \begin{alertblock}{#1}%
}{%
  \end{alertblock}%
}

% --- Speculative / unconfirmed claim. Research literacy matters here. ---
\newcommand{\rumour}{{\color{UHOrange}\textbf{[unconfirmed]}}\ }

% --- Key takeaway, for the closing slide of each section. ---
\newenvironment{takeaway}{%
  \begin{block}{Take away}%
}{%
  \end{block}%
}

% --- Live demo marker. ---
\newcommand{\demo}[1]{{\color{UHGreen}\textbf{Demo:}} #1}

% --- Emphasis consistent with the palette. ---
\newcommand{\key}[1]{{\color{UHBlue}\textbf{#1}}}
\newcommand{\warn}[1]{{\color{UHOrange}\textbf{#1}}}
\newcommand{\good}[1]{{\color{UHGreen}\textbf{#1}}}
\newcommand{\bad}[1]{{\color{UHRed}\textbf{#1}}}

% --- Units and quantities used constantly across the decks. ---
\newcommand{\unit}[1]{\,\mathrm{#1}}
\newcommand{\flops}{\mathrm{FLOPs}}
% NB: no \toks macro here --- \toks is a TeX primitive (a token register), so
% \newcommand{\toks} is a hard "already defined" error. No deck used it.
\newcommand{\KV}{\mathrm{KV}}
\newcommand{\softmax}{\operatorname{softmax}}

% =====================================================================
%  Code listings — used for config snippets, Slurm scripts, chat templates.
%  Frames containing these MUST be declared [fragile].
% =====================================================================
\lstdefinestyle{uh}{
  basicstyle=\ttfamily\scriptsize\color{UHDark},
  keywordstyle=\color{UHBlue}\bfseries,
  commentstyle=\color{UHGrey}\itshape,
  stringstyle=\color{UHGreen},
  numberstyle=\tiny\color{UHGrey},
  backgroundcolor=\color{UHLight},
  frame=none,
  xleftmargin=0.6em,
  framexleftmargin=0.6em,
  breaklines=true,
  breakatwhitespace=true,
  showstringspaces=false,
  columns=fullflexible,
  keepspaces=true,
  upquote=true,
  literate={-}{{-}}1 {_}{{\_}}1,
}
\lstset{style=uh}

\lstdefinelanguage{slurm}{
  morekeywords={SBATCH,module,srun,source,export,set,echo},
  morecomment=[l]{\#},
  sensitive=true,
}
\lstdefinelanguage{jinja}{
  morekeywords={for,endfor,if,endif,set,in},
  morecomment=[s]{\{\#}{\#\}},
  sensitive=true,
}

% =====================================================================
%  TikZ styles shared across decks (block diagrams, token grids).
% =====================================================================
\tikzset{
  mod/.style     = {draw=UHBlue, fill=UHLight, rounded corners=2pt,
                    minimum height=7mm, minimum width=22mm,
                    font=\scriptsize, align=center},
  modhi/.style   = {mod, draw=UHOrange, fill=UHOrange!12},
  modmute/.style = {mod, draw=UHGrey, fill=white, text=UHGrey},
  flow/.style    = {-{Latex[length=2mm]}, draw=UHGrey, thick},
  flowhi/.style  = {flow, draw=UHOrange},
  resid/.style   = {-{Latex[length=2mm]}, draw=UHBlue, thick, dashed},
  tokcell/.style = {draw=UHGrey!50, minimum size=4.2mm, font=\tiny, inner sep=0pt},
  tokon/.style   = {tokcell, fill=UHBlue!35},
  tokoff/.style  = {tokcell, fill=UHLight},
  tokmask/.style = {tokcell, fill=UHGrey!15, text=UHGrey!60},
  lbl/.style     = {font=\scriptsize\color{UHGrey}},
}

% =====================================================================
%  Presenter notes.
%
%  Notes are written on every slide. By default they are NOT shown, so
%  `latexmk -xelatex main.tex` produces a clean deck for projection.
%
%  To present with notes on a second screen (pdfpc, Skim, Impressive):
%      uncomment the line below, rebuild to main-notes.pdf
%  To print a notes handout:
%      \setbeameroption{show only notes}
% =====================================================================
% \setbeameroption{show notes on second screen=right}

\setbeamertemplate{note page}{%
  {\scriptsize\color{UHGrey}\coursecode\ \textbar\ \lecturename\par}%
  \vskip0.3em\hrule\vskip0.6em%
  {\small\insertnote}%
}

% ---------------------------------------------------------------- misc
\setbeamertemplate{caption}[numbered]
\setbeamertemplate{itemize items}[circle]
\setlength{\parskip}{0.1em}

% ---------------------------------------------------------------- density
% These decks are text-dense and a 16:9 frame is only 90mm tall, so the
% default vertical rhythm costs several lines per slide. Tightened globally
% here rather than slide by slide.
\setbeamersize{text margin left=0.7cm, text margin right=0.7cm}
\makeatletter
\def\@listi{\leftmargin\leftmargini
  \topsep 1.5\p@ \@plus1\p@ \@minus1\p@
  \parsep 0\p@
  \itemsep 1.5\p@ \@plus1\p@ \@minus1\p@}
\let\@listI\@listi
\makeatother

% The title page is a stack of short paragraphs, so a non-zero \parskip
% pushes it past the bottom of the frame. Neutralise it there only.
\addtobeamertemplate{title page}{\setlength{\parskip}{0pt}}{}

% Tighter lists — the decks are dense.
\setbeamertemplate{itemize/enumerate body begin}{\small}
\setbeamertemplate{itemize/enumerate subbody begin}{\footnotesize}

% Table helpers
\newcolumntype{L}[1]{>{\raggedright\arraybackslash}p{#1}}
\newcolumntype{C}[1]{>{\centering\arraybackslash}p{#1}}
\newcolumntype{R}[1]{>{\raggedleft\arraybackslash}p{#1}}

% Allow slightly overfull boxes without a warning storm
\hfuzz=2pt

% =====================================================================
%  Shared course-map frame.
%  Every lecture opens with this, with the current week highlighted:
%      \coursemapframe{1}   ... \coursemapframe{4}
%  Pass 0 to highlight nothing.
% =====================================================================
% \cmrow{current week}{this row's week}{dates}{Tuesday}{Thursday}{teacher}
% Highlights the row in UH blue bold when the two week numbers match.
% Deliberately avoids \rowcolor so that colortbl is not needed.
% The & and \\ must NOT sit inside the \ifnum branches: TeX scans alignment
% cells while skipping the false branch, so an & there breaks the conditional
% ("Incomplete \ifnum"). Keep the branches down to a macro name and let the
% separators live in the macro bodies, which expand after the test resolves.
\newcommand{\cmhi}[5]{%
  \textcolor{UHBlue}{\textbf{#1}} & \textcolor{UHBlue}{\textbf{#2}} &
  \textcolor{UHBlue}{\textbf{#3}} & \textcolor{UHBlue}{\textbf{#4}} &
  \textcolor{UHBlue}{\textbf{#5}}\\%
}
\newcommand{\cmplain}[5]{#1 & #2 & #3 & #4 & #5\\}
\newcommand{\cmrow}[6]{%
  \ifnum#1=#2\relax\expandafter\cmhi\else\expandafter\cmplain\fi
  {#2}{#3}{#4}{#5}{#6}%
}
\newcommand{\coursemapframe}[1]{%
\begin{frame}{Course overview}
\centering
\small
\setlength{\parskip}{0pt}
\renewcommand{\arraystretch}{1.08}
\begin{tabular}{@{}c L{1.7cm} L{3.1cm} L{3.7cm} L{1.4cm}@{}}
\toprule
\textbf{Wk} & \textbf{Dates} & \textbf{Tuesday} & \textbf{Thursday} & \textbf{Teacher}\\
\midrule
\cmrow{#1}{1}{27 / 29 Oct}{Language models}{Pretraining}{Aarne}
\cmrow{#1}{2}{3 / 5 Nov}{Post-training}{Lab 2: tune vs prompt}{Aarne}
\cmrow{#1}{3}{10 / 12 Nov}{Inference}{Lab 3: serve}{Aarne}
\cmrow{#1}{4}{17 / 19 Nov}{Evaluation}{Lab 4: evaluate}{Jussi}
\cmrow{#1}{5}{24 / 26 Nov}{RAG}{Lab 5: RAG}{Dmitry}
\cmrow{#1}{6}{1 / 3 Dec}{Agents}{Lab 6: agent}{Dmitry}
\cmrow{#1}{7}{8 / 10 Dec}{Course recap}{Special topics (TBC)}{All}
\bottomrule
\end{tabular}

\vspace{1.2em}
{\footnotesize Reusing a task and model across labs is recommended, not required.\\
There are no intermediate submissions; the final report and code repository
determine the grade.}
\note{Fifteen seconds, not more. The point of this slide is orientation, not
content: show them the shape of the course and where today sits inside it.
The final week is deliberately left flexible for a recap and special topics.}
\end{frame}%
}
